\begin{enumerate}
\item Xác định rõ phép thử: thực hiện mấy lần, có phân biệt thứ tự hay không, có hoàn lại hay không.
\item Viết $\Omega$ bằng cách liệt kê, dùng sơ đồ cây (phép thử nhiều giai đoạn) hoặc lập bảng (gieo $2$ xúc xắc).
\item Chọn trong $\Omega$ những kết quả thỏa mãn điều kiện để viết biến cố $A$; biến cố đối $\overline{A}$ gồm các kết quả còn lại.
\end{enumerate}
Lưu ý: tung $2$ đồng xu thì $SN$ và $NS$ là hai kết quả khác nhau. Gieo $2$ xúc xắc thì kết quả được ghi thành cặp có thứ tự $(i;j)$.
